% ---------------------------------------------------------------------------
% replication_report.tex
%
% The project's write-up. Skeleton at this stage: the structure and the
% build chain are in place and compile, but the discussion sections and our own
% exhibits are still to be written.
%
% Deliberately plain `article` rather than the scaffold's my_article_header.sty.
% That header currently fails to compile -- `\newtheorem{mycorollary}` in
% my_common_header.sty raises "Command \c@mylemma already defined" -- and this
% document needs none of the theorem machinery it provides. Worth revisiting if
% we later want the scaffold's styling.
%
% Every number and figure here is generated by code and pulled in with \input or
% \includegraphics. Nothing is typed by hand, so `doit` regenerates the whole
% document from the raw pulls.
%
% Build:  doit compile_latex_docs      (or: latexmk -xelatex -cd reports/replication_report.tex)
% ---------------------------------------------------------------------------
\documentclass[11pt]{article}

\usepackage[margin=1in]{geometry}
\usepackage{booktabs}
\usepackage{graphicx}
\usepackage{amsmath}
\usepackage{caption}
\usepackage[colorlinks=true, allcolors=blue]{hyperref}

% Generated artefacts live in _output/, resolved relative to this file so the
% build works from any working directory.
\newcommand{\PathToOutput}{../_output}
\graphicspath{{../_output/}}

\title{Predicting Returns with Financial Ratios:\\
       A Replication and Extension of Lewellen (2004)}
\author{Anthony Mazy \and Stefano Ramponi}
\date{\today}

\begin{document}
\maketitle

\begin{abstract}
We replicate Lewellen (2004), which argues that the standard Stambaugh (1999)
small-sample bias correction understates the predictive power of financial
ratios, and that imposing the economically motivated bound $\rho \le 1$ yields
substantially stronger evidence of return predictability. We reproduce the
paper's dividend-yield results on CRSP data, extend the sample past its 2000
cutoff, and report where our numbers diverge and why.
\end{abstract}

\section{Introduction}
% TODO: what the paper claims, why the bias correction matters, what we set out
% to reproduce.

\section{Data}
% TODO: CRSP (SIZ and CIZ schemas), Compustat via the CCM link, the
% trailing-12-month dividend construction, the 4-month accounting lag, and the
% book-equity definition question.

\section{Replication}
% TODO: Tables 1-3 reproduce; Tables 5/6 diverge on the level and why.

\section{Extension}
% TODO: 1946-present and 2001-present; the post-2000 collapse in $\hat\rho$ and
% what the paper's own Table A.1 predicts about it.

\section{Our Own Exhibits}
% TODO: this is where the summary table and figures go.

\subsection*{Placeholder: generated table}
The table below is produced by \texttt{src/pandas\_to\_latex\_demo.py} and
included with \verb|\input|. It exists at this stage only to prove that the
build chain works end to end -- code writes a \texttt{.tex} fragment into
\texttt{\_output/}, and this document pulls it in without anything being typed
by hand. It will be replaced by our own summary statistics.

\begin{table}[h]
\centering
\input{\PathToOutput/pandas_to_latex_simple_table1.tex}
\caption{Placeholder, to be replaced. Demonstrates that numbers reach this
document only by being generated.}
\end{table}

\section{Conclusion}
% TODO

\end{document}
